\section*{الفصل \ref{ch:g9:earth-resources} --- موارد من الأرض}

\begin{solution}{exo:g9:earth-resources:1}
(أ) \omterm{def:g9:earth-resources:aggregate}{الحصويات}. (ب) \omterm{def:g9:earth-resources:reserve}{الاحتياطي}. (ج) \omterm{def:g9:earth-resources:ore-deposit}{مكمن خام}. (د) \omterm{def:g9:earth-resources:ore-deposit}{معامل التركيز}.
\end{solution}

\begin{solution}{exo:g9:earth-resources:2}
(أ) في \omterm{def:g3:stones-pebbles-sand:pebble}{حصباء} \omterm{def:g1:landscapes-around-us:water}{الأنهار}. (ب) في الشقوق، من الماء الحار. (ج) في \omterm{def:g7:volcanoes:magma}{الصهارة}.
(د) هذه أول خطوة في تكوين مكمن في \omterm{def:g3:stones-pebbles-sand:pebble}{حصباء} \omterm{def:g1:landscapes-around-us:water}{نهر}: فالذهب المتحرر يحمله
\omterm{def:g1:landscapes-around-us:water}{النهر} بعد ذلك.
\end{solution}

\begin{solution}{exo:g9:earth-resources:3}
(ب)، (د)، (أ)، (ه)، (ج).
\end{solution}

\begin{solution}{exo:g9:earth-resources:4}
(أ) غير متجدد: يستغرق تكوينه ملايين السنين. (ب) متجدد: \omterm{def:g4:clouds-rain-snow:precipitation}{فالمطر} يعيد
ملء ما يؤخذ. (ج) غير متجدد: \omterm{def:g9:earth-resources:ore-deposit}{فمكمن الخام} يتكوّن على مدى أزمنة طويلة
جدًا. (د) غير متجدد بهذه الوتيرة: \omterm{def:g6:water-underground:water-table}{فمنسوب المياه الجوفية} ينخفض وتجف
\omterm{def:g6:water-underground:spring}{الآبار} في النهاية.
\end{solution}

\begin{solution}{exo:g9:earth-resources:5}
(أ) $4 \div 0.007 \approx 571$، أي نحو $10^3$ (بين $10^2$ و$10^3$،
وأقرب إلى $10^3$). (ب) $0.5 \div 0.0002 = 2500$، أي نحو $10^3$. (ج) $20
\div 5.6 \approx 3.6$، أي نحو 4: بين $10^0$ و$10^1$.
(د) الحديد وفير جدًا أصلًا في \omterm{def:g6:inside-the-earth:layers}{القشرة}: فيكفي خام أغنى من \omterm{def:g3:stones-pebbles-sand:rock}{الصخر} العادي
ببضع مرات فقط.
\end{solution}

\begin{solution}{exo:g9:earth-resources:6}
(أ) $19.3 \div 2.5 \approx 7.7$ مرة. (ب) الذهب: $19.3 \times 0.01 =
\qty{0.193}{g}$؛ \omterm{def:g3:stones-pebbles-sand:rock}{الصخر}: $2.5 \times 0.01 = \qty{0.025}{g}$. (ج) حين
يتحرك القاع كله، تغوص حبيبات الذهب الثقيلة بين \omterm{def:g3:stones-pebbles-sand:sand}{الرمل} \omterm{def:g3:stones-pebbles-sand:pebble}{والحصى} الأخف، حتى
\omterm{def:g3:stones-pebbles-sand:rock}{الصخر الأساسي} الصلب، حيث تتوقف.
\end{solution}

\begin{solution}{exo:g9:earth-resources:7}
(أ) الذهب والقصدير والكروم. (ب) بين $10^1$ و$10^2$، قريبًا من
$10^2$. (ج) $3000 \div 3.7 \approx 800$ مرة. (د) الكروم (أو القصدير، أو
الذهب): فكلما كبر \omterm{def:g9:earth-resources:ore-deposit}{معامل التركيز}، ندرت الأماكن التي ركّزت فيها الطبيعة
الفلز بما يكفي.
\end{solution}

\begin{solution}{exo:g9:earth-resources:8}
(أ) $18 \times 365 = \qty{6570}{kg}$، أي نحو \qty{6.6}{t}. (ب) $5 \times
10^{10}$ طن. (ج) $25000 \div 18 \approx 1389$ يومًا، أي قرابة أربع سنوات.
(د) \omterm{def:g3:stones-pebbles-sand:sand}{الرمل} المحفور من \omterm{def:g1:landscapes-around-us:water}{النهر} لم يعد يسافر مع التيار حتى \omterm{def:g1:landscapes-around-us:water}{الساحل}؛ والأمواج
تواصل حمل \omterm{def:g3:stones-pebbles-sand:sand}{الرمل} بعيدًا عن الشاطئ، لكن ما يصل ليعوّضه
أقل.
\end{solution}

\begin{solution}{exo:g9:earth-resources:9}
(أ) من نباتات وحيوانات بحرية دقيقة، مثل الدياتومات، دُفنت في \omterm{def:g3:stones-pebbles-sand:sand}{وحل} قاع
\omterm{def:g2:water-on-earth:ocean}{البحر}. (ب) بين نحو 60 و\qty{120}{\degreeCelsius}.
(ج) لأنهما أخف من الماء الذي يملأ \omterm{def:g6:water-underground:porosity}{مسام} \omterm{def:g3:stones-pebbles-sand:rock}{الصخور}.
(د) طبقة \omterm{def:g6:water-underground:permeability}{كتيمة}، مثل السجّيل؛ ويتجمعان تحتها، في قمة طية مقوّسة إلى
أعلى في كثير من الأحيان، والغاز فوق النفط.
\end{solution}

\begin{solution}{exo:g9:earth-resources:10}
(أ) $2.9 \times 10^7$ طن. (ب) $2.9 \times 10^7 \div 100 = 2.9 \times
10^5$ طن في السنة تقريبًا. (ج) الخام أثقل بمئة مرة من النحاس الذي
يحويه: $2.9 \times 10^7 \times 100 = 2.9 \times 10^9$ طن.
(د) كل طن من الخام لا يعطي إلا \qty{10}{kg} من النحاس، فيجب حفر
مليارات الأطنان من \omterm{def:g3:stones-pebbles-sand:rock}{الصخر}.
\end{solution}

\begin{solution}{exo:g9:earth-resources:11}
(أ) بمعدل 23~مليون طن في السنة كان سيكفي نحو $280 \div 23
\approx 12$ سنة، أي حتى مطلع الثمانينيات؛ ومع ذلك استُخرج 697~مليون طن،
أي أكثر بكثير من ذلك \omterm{def:g9:earth-resources:reserve}{الاحتياطي}. (ب) $980 \div 280 = 3.5$ مرة. (ج) $980 \div
23 \approx 43$ سنة. (د) مثلًا: تُكتشف مكامن جديدة؛ وتجعل التقنيات
الأفضل أو الأسعار الأعلى المكامنَ الأفقر جديرة بالتعدين.
\end{solution}

\begin{solution}{exo:g9:earth-resources:12}
(أ) $5 \times 10^1$ سنة. (ب) $2 \times 10^3 \div 50 = 40$ مرة. (ج) $3
\times 10^8 \div (5 \times 10^1) = 6 \times 10^6$ مرة. (د) استغرق
تكوين الفحم ملايين أضعاف المدة التي نستغرقها لاستهلاكه: فالطبيعة لا
تجدده على سلّمنا الزمني.
\end{solution}

\begin{solution}{pb:g9:earth-resources:1}
\textbf{1.} $100 \times 1000 = \num{100000}$~kg $= 1 \times 10^5$~kg.

\textbf{2.} $2000 \times 300 = \num{600000}$~t $= 6 \times 10^5$~t.

\textbf{3.} $6 \times 10^5 \div 10 = 6 \times 10^4$~t في السنة.

\textbf{4.} $0.40 \times 6 \times 10^4 = 2.4 \times 10^4$~t من \omterm{def:g3:stones-pebbles-sand:sand}{الرمل} في
السنة.

\textbf{5.} $1200 \times 500 \times 5 = 3 \times 10^6$~m$^3$.

\textbf{6.} ثلثاها: $2 \times 10^6$~m$^3$.

\textbf{7.} $2 \times 10^6 \times 1.6 = 3.2 \times 10^6$~t.

\textbf{8.} $\num{24000} + \num{136000} = \num{160000}$~t $= 1.6 \times
10^5$~t في السنة.

\textbf{9.} $\num{160000} \div 2000 = 80$: أي أكثر بثمانين مرة.

\textbf{10.} لا: فالحفرة تُفرغ أسرع بثمانين مرة مما يجدد \omterm{def:g1:landscapes-around-us:water}{النهر} \omterm{def:g3:stones-pebbles-sand:sand}{الرمل}.
فعلى سلّم المدينة، هذا \omterm{def:g3:stones-pebbles-sand:sand}{الرمل} مورد محدود.

\textbf{11.} $\num{24000} \div 4 = \num{6000}$~t من \omterm{def:g3:stones-pebbles-sand:sand}{الرمل} في السنة.

\textbf{12.} $3.2 \times 10^6 \div (1.6 \times 10^5) = 20$: \textbf{20 سنة}.
\end{solution}
